\subsection{How far the findings generalise}
\label{sec:scale}

Table~\ref{tab:general} repeats the measurement across model pairs, model families and test domains. The content contrast can grow because matched messages help more or because irrelevant messages hurt more, so for each bridge we report both its content ratio and its \emph{net} ratio: the total reduction in cross-entropy relative to no context, divided by the reference's. Appendix~\ref{app:general} lists the populations and protocol differences.

\begin{table}[t]
\centering\small
\caption{Generalisation. For each bridge, content gain (other-document donor minus matched message) and net gain (no context minus matched message), each relative to the receiver-state reference of the same receiver on the same windows; reference content in nats. 95\% intervals over documents, fitting seeds and donor draws. All bridges use the communication objective; layer pairs were selected on validation text for each model pair. The code and review rows evaluate the main pair's bridges, fitted on WikiText-like text, without refitting. Too few review windows have a same-document donor to split content into article-level and window-specific parts (Appendix~\ref{app:general}).}
\label{tab:general}
{\footnotesize\setlength{\tabcolsep}{4pt}\begin{tabular}{llccc}
\toprule
Setting & Probe text & Ref.\ content & Content ratio & Net ratio \\
\midrule
Pythia 70M$\to$160M (main) & Real text & 1.06 & 0.99 [0.94, 1.03] & 1.00 [0.97, 1.04] \\
 & LLM text &  & 0.82 [0.79, 0.85] & 0.72 [0.68, 0.76] \\
 & Self-samples &  & 0.44 [0.42, 0.47] & 0.31 [0.29, 0.33] \\
\addlinespace
Pythia 410M$\to$160M & Real text & 1.06 & 1.02 [0.98, 1.06] & 0.99 [0.96, 1.03] \\
 & LLM text &  & 0.89 [0.86, 0.92] & 0.70 [0.67, 0.73] \\
 & Self-samples &  & 0.55 [0.45, 0.65] & 0.31 [0.22, 0.40] \\
\addlinespace
Pythia 160M$\to$410M & Real text & 1.51 & 0.85 [0.79, 0.91] & 0.87 [0.82, 0.92] \\
 & LLM text &  & 0.77 [0.75, 0.80] & 0.69 [0.67, 0.72] \\
 & Self-samples &  & 0.27 [0.24, 0.34] & 0.14 [0.09, 0.17] \\
\addlinespace
GPT-2$\to$Pythia 160M & Real text & 1.05 & 1.02 [0.98, 1.05] & 1.00 [0.97, 1.03] \\
 & LLM text &  & 0.87 [0.78, 0.93] & 0.70 [0.61, 0.76] \\
 & Self-samples &  & 0.70 [0.66, 0.74] & 0.54 [0.51, 0.58] \\
\addlinespace
GPT-2$\to$Qwen2.5 0.5B & Real text & 1.58 & 0.50 [0.25, 0.76] & 0.62 [0.44, 0.80] \\
 & LLM text &  & 0.57 [0.46, 0.67] & 0.43 [0.34, 0.52] \\
\addlinespace
OPT 125M$\to$GPT-2 & Real text & 1.40 & 0.28 [0.26, 0.30] & 0.65 [0.60, 0.71] \\
 & LLM text &  & 0.28 [0.25, 0.31] & 0.27 [0.22, 0.32] \\
\addlinespace
Code (70M$\to$160M) & Real text & 1.26 & 0.31 [0.26, 0.34] & 0.03 [$-$0.03, 0.10] \\
 & LLM text &  & 0.30 [0.27, 0.32] & $-$0.10 [$-$0.16, $-$0.03] \\
 & Self-samples &  & 0.30 [0.28, 0.31] & 0.21 [0.19, 0.24] \\
\addlinespace
Reviews (70M$\to$160M) & Real text & 0.43 & 0.81 [0.76, 0.87] & 0.71 [0.65, 0.76] \\
 & LLM text &  & 0.73 [0.70, 0.76] & 0.63 [0.60, 0.67] \\
 & Self-samples &  & 0.46 [0.43, 0.50] & 0.46 [0.42, 0.50] \\
\bottomrule
\end{tabular}
}
\end{table}

\paragraph{The receiver's own states are matched only in some pairs.} In-domain, the real-text bridge's net ratio is close to 1 when pythia-70m, pythia-410m or GPT-2 sends to pythia-160m (1.00, 0.99 and 1.00). It is lower, 0.87 [0.82, 0.92], when pythia-160m sends to pythia-410m, a receiver whose own states are far more useful (1.51 nats of content, 93\% of what text gives it). The shortfall could come from the sender, the bridge's capacity or optimisation, or the restricted layer search, since the selected sender layer was the deepest candidate; we cannot tell these apart. The main pair shows that a less capable sender is not by itself an obstacle.

\paragraph{No net improvement was detected with a stronger sender.} We replaced pythia-70m by pythia-410m as sender to the same receiver, on identical windows and donors. The LLM-text bridge's content ratio rises from 0.82 to 0.89, but the net gain does not change detectably for any probe source, and its point estimates are slightly lower. For real text, LLM text and self-samples, the net gain changes by $-0.01$ [$-0.06$, 0.04], $-0.03$ [$-0.07$, 0.01] and $-0.00$ [$-0.09$, 0.11] nats, and the content gain by $+0.03$ [$-0.03$, 0.11], $+0.07$ [0.04, 0.11] and $+0.11$ [$-0.01$, 0.22]. The bootstrap pairs fitting seeds across senders, which share probe and minibatch seeds, and uses identical donor draws. Only the LLM-text bridge's content change is detectably positive, and in these fits its larger contrast reflects worse predictions with irrelevant messages rather than better ones with relevant messages. Read alone, the content contrast would suggest an improvement that the net gain does not show.

\paragraph{Other model families.} \XFAMILYRESULTS

\paragraph{Other domains.} The main pair's bridges were fitted for encyclopedia-like text. On long IMDB reviews (at least 322 tokens, 30\% of the test split), the domain's own reference is weaker, with 0.43 nats of content, and the bridges keep most of its content and net gain. Their content and net ratios are 0.81 and 0.71 for real text, 0.73 and 0.63 for LLM text, and 0.46 and 0.46 for self-samples.

On Python source code, all three retain about 30\% of the reference's content, but their net effects diverge. The self-sample bridge still helps (net ratio 0.21 [0.19, 0.24]), the real-text bridge's effect is unresolved (0.03 [$-0.03$, 0.10]), and the LLM-text bridge leaves the receiver worse off than no message ($-0.10$ [$-0.16$, $-0.03$]). On code, even an unrelated file's prefix given as text hurts the receiver (generic part $-0.26$). Relative performance falls sharply on this corpus, and which bridge is useful there is not predicted by which is best in-domain.

\paragraph{Text source and amount.} \TEXTSOURCES

\paragraph{Bridge capacity.} \CAPACITY
